\SourceLeaf{416}%
Geraden auch angegeben werden, in welchem Sinn man sie sich als durchlaufen denkt. Ohne diese Unterscheidung würde die Richtung einer Geraden stets zwei entgegengesetzten Radien entsprechen.

\Needspace{5\baselineskip}
\begin{center}2.\end{center}
\noindent\hspace*{1.2em}Im Hilfskreis nimmt man einen beliebigen Radius als den ersten oder seinen Endpunkt auf dem Umfang als Anfangspunkt an und entscheidet, in welchem Sinn man die Bögen von dort aus positiv zählen will (von links nach rechts oder umgekehrt). In entgegengesetzter Richtung gelten die Bögen also als negativ. So wird jede Richtung einer Geraden in Grad etc. oder durch eine Zahl ausgedrückt, die sie in Teilen des Radius angibt.

Angaben, die sich um $360^{0}$ oder ein Vielfaches davon unterscheiden, beziehen sich also eigentlich auf dieselben Richtungen und können im Allgemeinen als gleichwertig betrachtet werden. Wird jedoch auch berücksichtigt, wie eine veränderliche Neigung entsteht, kann es nötig sein, solche Angaben sorgfältig zu unterscheiden.

Hat man sich z. B. entschieden, die Bögen von links nach rechts zu zählen, und entsprechen den beiden Geraden $l,l'$ die Richtungen $L,L'$, so ist $L'-L$ die Neigung dieser Geraden zueinander. Man sieht leicht, dass, sofern $L'-L$ zwischen $-180^{0}$ und $+180^{0}$ liegt, der positive oder negative Wert zugleich anzeigt, ob $l'$ vom Schnittpunkt aus rechts oder links von $l$ liegt. Allgemein entscheidet dies das Vorzeichen von $\sin(L'-L)$.

Es sei $aa'$ ein Stück einer gekrümmten Linie, und den Tangenten in $a,a'$ entsprächen jeweils die Richtungen $\alpha,\alpha'$. Diese Buchstaben sollen die betreffenden Punkte des Hilfskreises bezeichnen, während ihre Bogenabstände vom Anfangspunkt mit $A,A'$ bezeichnet werden. Die Größe des Kreisbogens $\alpha\alpha'$ oder $A'-A$ heißt die \textls[100]{Amplitudo} von $aa'$.

Der Vergleich der Amplitudo des Bogens $aa'$ mit seiner Länge führt uns zum Begriff der Krümmung. Es sei $l$ ein beliebiger Punkt des\space%
